\section{Introduction}

When faced with repetition, programmers naturally turn to boilerplate or
parametrisation. In this paper, we focus on repetition that
arises from \emph{symmetry}. We develop
mathematical foundations that explicitly and uniformly capture
existing implicit patterns in programming and reasoning, and
define a construction that removes redundant symmetric cases from
functions and proofs.

Consider Gabriele Cirulli's famous game of 2048\footnote{Available to
play: \url{https://play2048.co}}, in which the player slides tiles
labelled by powers of $2$, merging tiles with identical labels. The
player can slide tiles in all four directions, and the game
source-code accommodates that ability through a general merging code
that takes the direction as a parameter, builds a sequence of matrix
indices to consider, and traverses over them\footnote{See
\texttt{GameManager.prototype.move}:
\url{https://github.com/gabrielecirulli/2048/blob/478b6ec/js/game_manager.js\#L129}.}.

Our anecdotal experience with half a dozen student projects
suggests this parametrisation can be difficult to implement
correctly. Here is an approach that is simpler to understand and to implement.
Define the merging logic only for a single direction, say left,
and exploit the inherent symmetry of the game for the other three
directions: rotate the board so that the task is to slide left,
perform the merge, and rotate back, restoring the original
orientation. Here is an example downwards merge:
\usetikzlibrary{arrows,arrows.meta}
\[
  \newcommand{\twofortyeight}[2]{\vcenter{\hbox{
        \begin{tikzpicture}[scale=0.6, tile/.style={rounded corners=0.05cm, minimum size=0.55cm, draw=white, thick, font=\sffamily}
          ]
          \foreach \x in {0,1,2,3} {
            \foreach \y in {0,1,2,3} {
              \node[tile, fill=gray!10] (\x\y) at (\x, \y) {};
            }
          }
          \foreach \row/\col/\num in {#1}
          \node[tile, fill=gray!10!orange!\num0] at (\col, 3 - \row) {\num};
          %%%%%%%%%%%%%%%%
          #2
  \end{tikzpicture}}}}
  \twofortyeight{0/1/2, 0/3/2, 1/3/2, 2/1/2, 2/3/4, 3/1/4, 3/3/8}
  {\draw[ double distance=1.5pt
      , arrows = {-Latex[]}
    ]  (02.center) node[above=-2pt]
    {
      \begin{tabular}{@{}>{\scriptsize\centering}p{.6cm}@{}}
        merge \\
        down
      \end{tabular}
    }
    -- (01.center)
    ;}
  \quad
  \explain\unRotation{rotate}
  \quad
  \twofortyeight{1/3/2, 3/3/2, 3/2/2, 1/1/2, 3/1/4, 1/0/4, 3/0/8}
  {\draw[ double distance=1.5pt
      , arrows = {-Latex[]}
  ]  (23.center) node[above] {\clap{\scriptsize merge left\qquad}} -- (13.center);}
  \quad
  \explain\mapsto{perform \\[-5pt] merge}
  \quad
  \twofortyeight{3/2/4, 1/1/4, 3/1/4, 1/0/4, 3/0/8}
  {}
  \quad
  \explain\Rotation{rotate \\[-5pt] back}
  \quad
  \twofortyeight{1/3/4, 2/1/4, 2/3/4, 3/1/4, 3/3/8}
  {}
\]

Another common example of symmetry is that of binary
search trees, in which elements smaller than the root go to the left
subtree, while larger elements go to the right. Operations such as
search, insertion, or deletion behave completely
symmetrically. One way of avoiding repetition by parametrisation is to
store both subtrees in an array of length two, and then use the
offset~$1 - i$ to refer to the tree opposite to the one at offset~$i$.
In practice, though, the additional complexity is often not worth
losing the high-level abstraction, and programmers prefer to duplicate
code.

But we can again avoid this repetition with the same approach and
implement the code only for one case. For example, when inserting an
element into the right subtree, we can flip the tree (and its ordering
relation!), insert the element into the left subtree, and then flip
the resulting tree back:
\usetikzlibrary{shapes.geometric, patterns, patterns.meta}
\[
  \newcommand{\insertion}[1]{\vcenter{\hbox{
        \begin{tikzpicture}[
            scale=0.75,
            inner/.style={circle, draw, minimum width=2.5mm},
            tree/.style={isosceles triangle, isosceles triangle stretches=true, draw, shape border rotate=90,minimum width=10mm},
            h/.style={minimum height=10mm},
            h1/.style={minimum height=12.5mm, yshift=-2.5mm},
            level 1/.style={sibling distance=20mm},
            transform/.style={bend left, dashed, shorten >=2.5mm, shorten <=2.5mm, -{stealth}}
          ] #1
  \end{tikzpicture}}}}
  \insertion{
    \node[inner] (t3) at (0, -4) {}
    [child anchor=north]
    child {node[pattern={Lines[angle=-45, distance=1pt]}, pattern color=red, tree, h] {}}
    child {node[pattern={Lines[angle=-45]}, pattern color=blue, tree, h] {}};
    \node (x') at (0.6, -4) {${} < x$};
    \path (x') edge[->, out=0, in=45] (1.3, -5);
  }
  \;\;\:\,
  \explain{\curvearrowright\hskip-10pt|}{flip\\[-5pt]everything}
  \;\;\:\,
  \insertion{\node[inner] (t1) at (0, 0) {}
    [child anchor=north]
    child {node[pattern={Lines[angle=45]}, pattern color=blue, tree, h] {}}
    child {node[pattern={Lines[angle=45, distance=1pt]}, pattern color=red, tree, h] {}};
    \node (x) at (-0.8, 0) {$x <^{\mathrm{op}} {}$};
    \path (x) edge[->, out=180, in=135] (-1.3, -1);
  }
  \;\;\:\,
  \explain\mapsto{insert\\[-5pt]w.r.t.~$(<^{\mathrm{op}})$}
  \;\;\:\,
  \insertion{
    \node[inner] (t2) at (4, 0) {}
    [child anchor=north]
    child {node[pattern={Lines[angle=45]}, pattern color=blue, tree, h1] {}}
    child {node[pattern={Lines[angle=45, distance=1pt]}, pattern color=red, tree, h] {}};
  }
  \;\;\:\,
  \explain{\curvearrowleft\hskip-8.5pt|}{flip\\[-5pt]back}
  \;\;\:\,
  \insertion{
    \node[inner] (t4) at (4, -4) {}
    [child anchor=north]
    child {node[pattern={Lines[angle=-45, distance=1pt]}, pattern color=red, tree, h] {}}
    child {node[pattern={Lines[angle=-45]}, pattern color=blue, tree, h1] {}};
  }
\]
A reader may justifiably worry about the cost of flipping the whole
tree each time we descend right, but a more efficient implementation
fits the same pattern.

Our final example comes from mathematics,
where one often invokes a \emph{without loss of generality} (or \WLOG)
argument when we treat a canonical case, and derive the remaining
cases by symmetry. Consider the pigeonhole principle given as an
example in the Wikipedia entry for \WLOG{}~\cite{wiki:wlog}:
%
\begin{quote}
If three objects are each painted either red or blue, then there must
be at least two objects of the same color.

\textbf{Proof.} Assume, without loss of generality, that the first
object is red. If either of the other two objects is red, then we are
finished; if not, then the other two objects must both be blue and we
are also finished.
\end{quote}
%
On paper, \WLOG appeals to an intuitive understanding that nothing has
been lost by considering a special case. But if the symmetry is complex
or if we want to convince a proof assistant, we need
to explicate the reduction of the general case
to the special one.

We can generalise the special case as a parametric proof that covers
all cases. Above, we could take the colour of the first object as a
parameter and then instantiate it to both red and blue. But as
\citeauthor{DBLP:conf/tphol/Harrison09} argues in his seminal study of
\WLOG in the context of the HOL proof assistant---and we agree---this
formulation is not faithful to the informal
proof~\cite{DBLP:conf/tphol/Harrison09}.

Modern proof assistants offer special tactics to treat \WLOG
arguments, but approach them from various directions. Some tactics or
theorems named \texttt{wlog} work in a very restricted setting, such
as symmetric relations or negation, and some cut in arbitrary
additional statements that do not involve symmetries.
These divergent approaches point towards a missing
mathematical foundation. Furthermore, none treat the
underlying symmetry as a first-class citizen, and as
\citeauthor{dijkstra:wlog} argues---and we agree---``omissions [of the
underlying symmetry] make most `\WLOG{} proofs' rather
unsatisfactory''~\cite{dijkstra:wlog}.

We can explicate this symmetry by following the same
pattern we used for programming. In the excluded case where the first object is blue, we can
exchange the colours to make it red, use the specialised argument, and
then revert the colours back to the original. For example, one of
the three cases where there is another blue object in addition to the
first one proceeds as:
\[
  \newcommand{\balls}[2]{\vcenter{\hbox{
    \begin{tikzpicture}[
      scale=0.5,
      obj/.style={circle, draw, inner sep=0pt, minimum size=4mm},
      rdot/.style={obj, fill=red!60, draw=red},
      bdot/.style={obj, fill=blue!20, draw=blue},
      pair/.style={rounded corners, thin, draw=gray},
      >=stealth, font=\small
    ]
      \pgfmathsetmacro\dx{-0.03}
      \foreach \col/\f in {#1}
        \node[\f] at (\col+0.5, 0.5) {\ifthenelse{\equal{\f}{bdot}}{b}{r}};
      #2
    \end{tikzpicture}
  }}}
  \balls{0/bdot, 1/rdot, 2/bdot}{}
  \qquad
  \explain{\rightleftharpoons}{exchange\\[-5pt]colours}
  \qquad
  \balls{0/rdot, 1/bdot, 2/rdot}{}
  \qquad
  \explain{\mapsto}{find another red\\[-5pt]object, if it exists}
  \qquad
  \balls{0/rdot, 1/bdot, 2/rdot}{
    \draw[pair] (0-\dx,0-\dx) rectangle (1+\dx,1+\dx);
    \draw[pair] (2-\dx,0-\dx) rectangle (3+\dx,1+\dx);
  }
  \qquad
  \explain{\leftrightharpoons}{exchange\\[-5pt]colours back}
  \qquad
  \balls{0/bdot, 1/rdot, 2/bdot}{
    \draw[pair] (0-\dx,0-\dx) rectangle (1+\dx,1+\dx);
    \draw[pair] (2-\dx,0-\dx) rectangle (3+\dx,1+\dx);
  }
\]
while if both remaining objects are red, we similarly proceed as:
\[
  \newcommand{\balls}[2]{\vcenter{\hbox{
    \begin{tikzpicture}[
      scale=0.5,
      obj/.style={circle, draw, inner sep=0pt, minimum size=4mm},
      rdot/.style={obj, fill=red!60, draw=red},
      bdot/.style={obj, fill=blue!20, draw=blue},
      pair/.style={rounded corners, thin, draw=gray},
      >=stealth, font=\small
    ]
      \pgfmathsetmacro\dx{-0.03}
      \foreach \col/\f in {#1}
        \node[\f] at (\col+0.5, 0.5) {\ifthenelse{\equal{\f}{bdot}}{b}{r}};
      #2
    \end{tikzpicture}
  }}}
  \balls{0/bdot, 1/rdot, 2/rdot}{}
  \qquad
  \rightleftharpoons
  \qquad
  \balls{0/rdot, 1/bdot, 2/bdot}{}
  \qquad
  \mapsto
  \qquad
  \balls{0/rdot, 1/bdot, 2/bdot}{
    \draw[pair] (1-\dx,0-\dx) rectangle (3+\dx,1+\dx);
  }
  \qquad
  \leftrightharpoons
  \qquad
  \balls{0/bdot, 1/rdot, 2/rdot}{
    \draw[pair] (1-\dx,0-\dx) rectangle (3+\dx,1+\dx);
  }
\]

In this paper, we develop this pattern into a uniform construction, and
use it to eliminate symmetric cases in definitions of functions and in reasoning.
On the surface, this pattern is conceptually simple, and so one may,
just like us, be quickly tempted to implement it, either in a
programming language or in a proof assistant. We
experimented in simply- and dependently-typed languages, and
unsurprisingly, each obstacle pushed us back to pen-and-paper redesigns and
forced us to enrich the underlying mathematical foundation.

For this reason, while still keeping computational applications in
mind, this work focuses on these mathematical foundations.
We postpone a bona fide implementation to future work. Such future
implementations will allow us to explore symmetric programming on more
involved examples, in a machine-checked environment, design the best
interface for exposing it to the programmer, to build a comprehensive
library for symmetric programming and mechanised reasoning, and
revisit existing \WLOG formalisations systematically. While we keep these
downstream goals in mind, we pursue none of them in the
present work, and remain solely mathematical.

Our contributions:
\begin{itemize}
  \item In \secref{wlog construction}, we identify three ingredients:
    \emph{group actions} capturing the symmetries of the input and
    output, a \emph{canoniser} map that assigns a canonising group
    element to each input, and a \emph{core function} that defines the
    construction on canonical inputs. We then characterise our symmetric \WLOG{}
    construction as the unique symmetry-preserving
    (i.e. equivariant) extension of the core function to all the inputs.
  \item In \secref{dependent types}, we extend this theory to the
    setting of \emph{dependent types}. We show how via
    the Curry-Howard correspondence, symmetry-invariant propositions
    correspond to types equipped with a group action. We show how proofs
    via the symmetric \WLOG{} argument amount to dependent functions obtained by
    the above construction.
    % % Sean Moss tells me this is probably known:
    % We also rephrase
    % the classical Orbit-Stabiliser theorem as an equivariant
    % bijection.
  \item In \secref{algebraic data types}, we further extend the theory
    to (simply- and dependently-typed) \emph{algebraic data types} and
    their induction principles. The main construct we introduce is
    \emph{symmetry strength}, which defines a group action in terms of
    data constructors, paving the way for symmetric reasoning on data types,
    and efficient partial evaluation.
\end{itemize}
%
We conclude with related and future work in \secref{conclusion}.
\ifthenelse{\boolean{WithSupplementary}}{%
  The Appendix has further details and proofs.%
}{%
  The accompanying Technical Report~\cite{kammar-pretnar:symmetric-programming-and-reasoning-tech-report}
  is an otherwise identical copy of this manuscript, apart from
  this sentence and detailed appendices supporting the claims we make, which we cross-reference from here.%
}


\begin{remark}
  Our development has been informed by category theoretic
  structure. To keep exposition accessible, we do not rely on the
  language of and concepts from category theory in the lion's share of this
  work.  The notable exception is our account of symmetric data types
  (\secref{algebraic data types}). There, generic accounts of the
  computational structure are close to their underlying categorical
  structure, and consequently we explicitly employ the categorical
  language. Readers who wish to reconstruct the category theoretic
  picture can consult \figref{indexed-datatypes-summary} in
  \subsecref{indexed data types} alongside the development.
\end{remark}
